\subsection{PRIMITIVES OF DIFFERENTIAL FORMS WITH VALUES IN A LIE ALGEBRA}

Lemma 5. Let X be a manifold of class \(\mathbf{C}^r\) , F and \(\mathbf{F}'\) vector bundles of class \(\mathbf{C}^r\) with base space X and \(\phi\) a morphism of F into \(\mathbf{F}'\) . For all \(x \in \mathbf{X}\) , let \(\mathbf{S}_x\) be the set of

\[
(a, \phi (a)) \in \mathrm{F} _ {x} \oplus \mathrm{F} _ {x} ^ {\prime}
\]

for \(a \in \mathbf{F}_x\) . Then the union \(\mathbf{S}\) of the \(\mathbf{S}_x\) is a vector subbundle of \(\mathbf{F} \oplus \mathbf{F}'\) .

Let \(\theta\) and \(\theta'\) be the mappings of \(F \oplus F'\) into itself defined as follows: if \((u, v) \in F_x \oplus F_x'\) , then

\[
\theta (u, v) = (u, v + \phi (u)), \quad \theta^ {\prime} (u, v) = (u, v - \phi (u)).
\]

By Differentiable and Analytic Manifolds, R, 7.7.1, \(\theta\) and \(\theta'\) are morphisms of \(F \oplus F'\) into itself. Clearly \(\theta \circ \theta' = \theta' \circ \theta = \mathrm{Id}_{F\oplus F'}\) . Hence \(\theta\) and \(\theta'\) are automorphisms of \(F \oplus F'\) . Therefore, \(S = \theta(F \oplus \{0\})\) is a vector subbundle of \(F \oplus F'\) .

Lemma 6. Let G be a Lie group germ, \(\omega\) the canonical left differential form of G (§ 3, no. 18.9), M a manifold of class \(\mathbf{C}^r\) ( \(r \geqslant 2\) ) and \(\alpha\) a differential form of class \(\mathbf{C}^{r-1}\) and degree 1 on M with values in L(G).

\begin{enumerate}
\def\labelenumi{(\roman{enumi})}
\tightlist
\item
  The elements of \(\mathbf{T}(\mathbf{M}\times \mathbf{G})\) at which the differential form
\end{enumerate}

\[
\theta = \mathrm{pr} _ {1} ^ {*} \alpha - \mathrm{pr} _ {2} ^ {*} \omega
\]

is zero constitute a vector subbundle S of T(M × G) of class \(\mathbf{C}^{r-1}\) .

\begin{enumerate}
\def\labelenumi{(\roman{enumi})}
\setcounter{enumi}{1}
\item
  For all \((x,g)\in \mathbf{M}\times \mathbf{G},\mathrm{T}(\mathrm{pr}_1)|\mathrm{S}_{(x,g)}\) is an isomorphism of \(\mathrm{S}_{(x,g)}\) onto \(\mathrm{T}_x(\mathbf{M})\) .
\item
  If \(d\alpha + [\alpha]^2 = 0\) (cf.~§3, no. 14) the vector subbundle S is integrable.
\end{enumerate}

If \((x,g)\in \mathbf{M}\times \mathbf{G}\) and \((u,v)\in \mathrm{T}_x(\mathrm{M})\times \mathrm{T}_g(\mathrm{G})\) , then

\[
\theta_ {(x, g)} (u, v) = \alpha (u) - g ^ {- 1} v.
\]

Hence the kernel of \(\theta_{(x,g)}\) is the set \(\mathrm{S}_{(x,g)}\) of \((u,g\alpha (u))\) for \(u\in \mathrm{T}_x(\mathbf{M})\) , whence (ii). We consider \(\mathrm{T}(\mathrm{M}\times \mathrm{G})\) as the direct sum of two vector bundles F and \(\mathbf{F}'\) with \(\mathrm{F}_{(x,g)} = \mathrm{T}_x(\mathrm{M})\times \{0\}\) and \(\mathrm{F}_{(x,g)}^{\prime} = \{0\} \times \mathrm{T}_g(\mathrm{G})\) for all

\[
(x, g) \in \mathbf {M} \times \mathbf {G}.
\]

For \(u \in \mathrm{T}_{x}(\mathbf{M}) \times \{0\}\) , we write \(\phi(u) = (0, g\alpha(u))\) . Using the left trivialization of \(\mathrm{T}(\mathrm{G})\) , it is seen that \(\phi\) is a morphism of \(\mathbf{F}\) into \(\mathbf{F}'\) , whence (i) (Lemma 5). Finally, if \(d\alpha + [\alpha]^2 = 0\) , then

\[
\begin{array}{r l} {d \theta = \mathrm{pr} _ {1} ^ {*} (d \alpha) - \mathrm{pr} _ {2} ^ {*} (d \omega)} \\ & {= - \frac {1}{2} (\mathrm{pr} _ {1} ^ {*} \alpha \wedge \mathrm{pr} _ {1} ^ {*} \alpha - \mathrm{pr} _ {2} ^ {*} \omega \wedge \mathrm{pr} _ {2} ^ {*} \omega)} \\ & {= - \frac {1}{2} (\mathrm{pr} _ {1} ^ {*} \alpha - \mathrm{pr} _ {2} ^ {*} \omega) \wedge (\mathrm{pr} _ {1} ^ {*} \alpha + \mathrm{pr} _ {2} ^ {*} \omega)} \\ & {= - \frac {1}{2} \theta \wedge (\mathrm{pr} _ {1} ^ {*} \alpha + \mathrm{pr} _ {2} ^ {*} \omega)} \end{array}
\]

and hence S is integrable (Differentiable and Analytic Manifolds, R, 9.3.6).

THEOREM 5. Let \(G\) be a Lie group germ, \(M\) a manifold of class \(C^r (r\geqslant 2)\) and \(\alpha\) a differential form of class \(\mathbf{C}^{r-1}\) and degree 1 on \(M\) with values in \(L(G)\) , such that \(d\alpha +[\alpha ]^2 = 0\) . For all \(x\in M\) and all \(g\in G\) , there exists a mapping \(f\) , defined and of class \(\mathbf{C}^{r-1}\) on an open neighbourhood of \(x\) , with values in \(G\) , such that \(f(x) = g\) and \(f^{-1}.df = \alpha\) . Two mappings which satisfy these conditions coincide on a neighbourhood of \(x\) .

Let \(x \in M\) and \(g \in G\) . By Lemma 6 (whose notation we adopt) and Differentiable and Analytic Manifolds, R, 9.3.7, there exist an open neighbourhood U of \(x\) in M and a mapping \(m \mapsto \phi(m) = (m, f(m))\) of class \(C^{r-1}\) of U into \(M \times G\) such that \(f(x) = g\) and \(\phi^*(\theta) = 0\) . Then

\[
\begin{array}{l l} f ^ {- 1}. d f = f ^ {*} (\omega) & \text {(§ 3, no. 18.9)} \\ = (\mathrm{pr} _ {2} \circ \phi) ^ {*} (\omega) & \text {(for \(f = \mathrm{pr} _ {2} \circ \phi\))} \\ = \phi^ {*} (\mathrm{pr} _ {1} ^ {*} \alpha - \theta) & \text {(Lemma 6)} \\ = \phi^ {*} (\mathrm{pr} _ {1} ^ {*} \alpha) & \text {(for \(\phi^ {*} (\theta) = 0\))} \\ = \alpha & \text {(for \(\mathrm{pr} _ {1} \circ \phi = \mathrm{Id} _ {\mathbf {U}}\))}. \end{array}
\]

Let \(f'\) be a mapping of class \(\mathbf{C}^{r-1}\) of U into G such that \(f'(x) = g\) and \(f'^{-1}df' = \alpha\) . By § 3, 18.9, \(ff'^{-1}\) is locally constant and hence \(f' = f\) in a neighbourhood of \(x\) .

PROPOSITION 10. Let M be an analytic manifold, g a complete normable Lie algebra and \(\alpha\) an analytic differential form of degree 1 on M, with values in g, with the following properties:

\begin{enumerate}
\def\labelenumi{(\alph{enumi})}
\item
  for all \(m \in \mathbf{M}\) , \(\alpha_{m}\) is an isomorphism of \(\mathrm{T}_{m}(\mathrm{M})\) onto \(\mathfrak{g}\) ;
\item
  \(d\alpha + [\alpha]^{2} = 0.\)
\end{enumerate}

Then, for all \(m_0 \in \mathbf{M}\) , there exist an open neighbourhood \(\mathbf{M}'\) of \(m_0\) in \(\mathbf{M}\) and a Lie group germ structure on \(\mathbf{M}'\) , compatible with the manifold structure on \(\mathbf{M}'\) , with identity element \(m_0\) and with the following properties:

\begin{enumerate}
\def\labelenumi{(\roman{enumi})}
\item
  \(\alpha_{m_{0}}\) is an isomorphism of L(M') onto g;
\item
  the differential form \(m \mapsto \alpha_{m_0}^{-1} \circ \alpha_m\) is the left canonical differential form of \(M'\) . If \(M_1'\) and \(M_2'\) are two such group germs, \(M_1'\) and \(M_2'\) have a common open subgroup germ.
\end{enumerate}

There exists a Lie group germ G such that \(L(G) = g\) . Let \(m_{0} \in M\) . By Theorem 5, there exist an open neighbourhood \(M'\) of \(m_{0}\) in M and an analytic mapping f of \(M'\) into G such that \(f(m_{0}) = e\) and \(f^{-1}.df = \alpha\) . Then \(T_{m_{0}}(f) = \alpha_{m_{0}}\) is an isomorphism of \(T_{m_{0}}(M)\) onto g; hence, shrinking \(M'\) and G, it can be assumed that f is an isomorphism of the manifold \(M'\) onto the manifold G. We transport to \(M'\) the Lie group germ structure on G by means of \(f^{-1}\) . Then \(T_{m_{0}}(f)\) becomes an isomorphism of \(L(M')\) onto \(L(G) = g\) , whence (i). On the other hand, if \(\omega\) denotes the left canonical differential form of G, then

\[
\begin{array}{r} \alpha_ {m _ {0}} ^ {- 1} \circ \alpha_ {m} = (\mathrm{T} _ {m _ {0}} f) ^ {- 1} \circ (f ^ {- 1}. d f) (m) \\ = (\mathrm{T} _ {m} f) ^ {- 1} \circ \omega (f (m)) \circ \mathrm{T} _ {m} f \end{array}
\]

and hence \(m \mapsto \alpha_{m_0}^{-1} \circ \alpha_m\) is the left canonical differential form of \(\mathbf{M}'\) .

Let \(\mathbf{M}^{\prime}\) be an open neighbourhood of \(m_0\) , with a Lie group germ structure, with identity element \(m_0\) and with the analogous properties to properties (i) and (ii). Then \(\alpha_{m_0}\) is an isomorphism of \(\mathrm{L}(\mathbf{M}')\) onto \(\mathfrak{g}\) and also of \(\mathrm{L}(\mathbf{M}'')\) onto \(\mathfrak{g}\) and hence \(\mathrm{L}(\mathbf{M}') = \mathrm{L}(\mathbf{M}'')\) . Therefore, shrinking \(\mathbf{M}'\) and \(\mathbf{M}''\) , it can be assumed that there exists an isomorphism \(\phi\) of the group germ \(\mathbf{M}'\) onto the group germ \(\mathbf{M}''\) (no. 1, Corollary 1 to Theorem 1). Then \(\phi^{-1}.d\phi\) is the canonical left differential of \(\mathbf{M}'\) . On the other hand, let \(\psi\) be the canonical injection of the manifold \(\mathbf{M}' \cap \mathbf{M}''\) into the Lie group germ \(\mathbf{M}''\) ; clearly \(\psi^{-1}.d\psi\) is a restriction of the canonical left differential of \(\mathbf{M}''\) . Hence \((\psi^{-1}.d\psi)(m) = \alpha_{m_0}^{-1} \circ \alpha_m = (\phi^{-1}.d\phi)(m)\) for all \(m \in \mathbf{M}' \cap \mathbf{M}''\) . Therefore \(\phi\) and \(\psi\) coincide on a neighbourhood of \(m_0\) (§ 3, 18.9). This proves the last assertion of the proposition.

COROLLARY. Let M be an analytic manifold of finite dimension n.~Let \(\omega_{1},\ldots,\omega_{n}\) be analytic differential forms of degree 1 on M, with scalar values, which are linearly independent at each point of M and such that, for all \(k=1,\ldots,n\) , \(\omega_{k}\) is a linear combination with constant coefficients of the \(\omega_{i}\wedge\omega_{j}\) . Then, for all \(m_{0}\in M\) , there exists an open neighbourhood \(M'\) of \(m_{0}\) in M and a Lie group germ structure on \(M'\) compatible with the manifold structure on \(M'\) , with identity element \(m_{0}\) and such that \(\omega_{1}|_{M'},\ldots,\omega_{n}|_{M'}\) form a basis of the space of left invariant differential forms on \(M'\) of degree 1 with scalar values.

If \(\mathbf{M}_1'\) and \(\mathbf{M}_2'\) are two such group germs, \(\mathbf{M}_1'\) and \(\mathbf{M}_2'\) have a common open subgroup germ.

Let \(\mathbf{X}_1, \ldots, \mathbf{X}_n\) be the vector fields on M such that, at each point \(m\) of M, the \((\mathbf{X}_i)_m\) constitute the basis of \(\mathrm{T}_m(\mathrm{M})\) dual to \(((\omega_1)_m), \ldots, (\omega_n)_m)\) . These fields are analytic. By hypothesis, there exist \(c_{ijk} \in \mathbf{K}\) ( \(1 \leqslant i,j, k \leqslant n\) ) such that \(c_{ijk} = -c_{jik}\) and \(d\omega_k = \sum_{i < j} c_{ijk}\omega_i \wedge \omega_j\) . By Differentiable and Analytic Manifolds, R, 8.5.7, formula (11),

\[
\langle [ \mathrm{X} _ {i}, \mathrm{X} _ {j} ], \omega_ {k} \rangle = - (d \omega_ {k}) (\mathrm{X} _ {i}, \mathrm{X} _ {j}) = - \left(\sum_ {r <   s} c _ {r s k} \omega_ {r} \wedge \omega_ {s}\right) (\mathrm{X} _ {i}, \mathrm{X} _ {j}) = - c _ {i j k}
\]

and hence \([\mathrm{X}_i,\mathrm{X}_j] = -\sum_{k}c_{ijk}\mathrm{X}_k\) . It then follows that the \(-c_{ijk}\) are the constants of structure of a Lie algebra \(\mathfrak{g}\) relative to a basis \((e_1,\dots ,e_n)\) . For all \(m\in \mathbf{M}\) , let \(\alpha_{m}\) be the linear mapping of \(\mathrm{T}_{m}(\mathrm{M})\) into \(\mathfrak{g}\) which maps \((\mathrm{X}_1)_m\) to \(e_1,\ldots ,(\mathrm{X}_n)_m\) to \(e_n\) . Then \(\alpha\) is an analytic differential form on \(\mathbf{M}\) of degree

1 with values in \(\mathfrak{g}\) and \(\alpha_{m}\) is an isomorphism of \(\mathrm{T}_m(\mathbf{M})\) onto \(\mathfrak{g}\) . On the other hand, \(\alpha = \sum_{k=1}^{n} \omega_k e_k\) and hence

\[
d \alpha = \sum_ {k = 1} ^ {n} (d \omega_ {k}) e _ {k} = \sum_ {k = 1} ^ {n} \left(\sum_ {i <   j} c _ {i j k} \omega_ {i} \wedge \omega_ {j}\right) e _ {k}
\]

and

\[
\begin{array}{r l} {[ \alpha ] ^ {2} = \sum_ {k = 1} ^ {n} [ \omega_ {k} e _ {k} ] ^ {2} + \sum_ {i <   j} (\omega_ {i} e _ {i}) \wedge (\omega_ {j} e _ {j})} & {\text {(§ 3, formula (30))}} \\ {= \sum_ {i <   j} (\omega_ {i} \wedge \omega _ {j}) [ e _ {i}, e _ {j} ]} \\ {= - \sum_ {k = 1} ^ {n} \sum_ {i <   j} (c _ {i j k} \omega_ {i} \wedge \omega_ {j}) e _ {k}} \\ {= - d \alpha.} \end{array}\tag{5}
\]

It then suffices to apply Proposition 10.

